
% File: executive_summary.tex

\newpage
\thispagestyle{plain}

\addcontentsline{toc}{section}{Executive Summary}

\begin{center}
\textbf{\LARGE Executive Summary}
\end{center}

\vspace{1cm}

\begin{spacing}{1.5}

Traditional student feedback collection methods in educational institutions rely heavily on manual paper-based forms or generic online surveys. These approaches suffer from low response rates, delayed processing, lack of analytical depth, and an inability to provide timely, actionable insights to educators. As a result, teachers often receive feedback that is either too late to act upon or too unstructured to derive meaningful conclusions.

The TechAware AI-Based Feedback System addresses these challenges by providing a comprehensive, web-based platform that automates the entire feedback lifecycle — from collection to analysis to actionable reporting. The system integrates Natural Language Processing (NLP) for real-time sentiment analysis, enabling automatic classification of student feedback as Positive (Understood), Negative (Not Understood), or Neutral. Based on the analysis, AI-powered suggestions are generated to help teachers improve their teaching methodologies.

The system supports four distinct user roles: Students, Teachers, Administrators, and Batch Advisors. Students submit structured feedback after lectures, which includes star ratings across multiple dimensions (clarity, punctuality, material quality, communication, and overall satisfaction) along with free-text comments. Teachers access a dedicated dashboard to view aggregated analytics, sentiment distributions, and AI-generated improvement suggestions. Administrators manage the entire academic infrastructure, including departments, batches, sections, courses, timetables, and user accounts. Batch Advisors receive automated alerts when teacher performance falls below defined thresholds.

The system is built using Python Flask as the backend framework, SQLAlchemy ORM with SQLite for data persistence, TextBlob for sentiment analysis, and HTML/CSS/JavaScript for the frontend. Additional integrations include Chart.js for data visualization, ReportLab for PDF report generation, and SMTP for automated email notifications.

The expected impact of this system is to enhance teaching quality through data-driven feedback, reduce administrative overhead in feedback management, and promote a culture of continuous improvement in academic institutions.

\end{spacing}
\newpage
